\section{Einleitung}

Digitale Signalverarbeitung (DSV) ermöglicht die softwarebasierte Verarbeitung
und Erzeugung von Signalen und findet unter anderem in der Audioverarbeitung 
Anwendung. Durch die zunehmende Leistungsfähigkeit von Mikrocontrollern ist
es möglich, solche Signalverarbeitungsverfahren auch auf eingebetteten Systemen 
umzusetzen. 

Ziel dieser Arbeit ist die Entwicklung einer einfach nutzbaren und erweiterbaren
DSV-Bibliothek für Mikrocontroller. In Bezug auf
Nutzerfreundlichkeit ist Java Streams \cite{java2026stream} ein 
sehr gutes Vorbild. Java Streams setzen als Fluent API \cite{fluentAPI} die Mentalität 
\glqq Signal geht rein, Signal geht raus\grqq\ (oder im fall von Java Streams
Collection) und die sequentielle Abarbeitung von Streams perfekt um.
Da C++ auch objektorientierte Programmierung unterstützt, ist es möglich, eine
ähnliche Fluent API in C++ zu implementieren. 
In den Abbildung \ref{Vergleich_Streams_Signal} wird diese sequentielle
Abarbeitung dargestellt und gezeigt wie eine Anwendung der DSV-Bibliothek 
aussehen könnte.

\begin{figure}[h]
  
  \begin{minipage}[h]{0.48\textwidth}
    \begin{lstlisting}[language=Java, 
      frame=single,
      captionpos=b]
      Stream.of(1,2,3,4)
        .map(x -> x*x)
        .filter(x -> x<10)
        .toArray();
    \end{lstlisting}
  \end{minipage}
  \hfill
  \begin{minipage}[h]{0.48\textwidth}
    \begin{lstlisting}[language=C++, 
      frame=single,
      captionpos=b]
      Signal.fromInput(ADC)
        .multiply(2)
        .filter(20,5000)
        .applyFunction(coolFkt)
        .toOutput(DAC);
    \end{lstlisting}
  \end{minipage}
  \caption{Vergleich Einer Beispielanverndung von Javastreams und 
            der zu entwickelnden Signalbibliothek}
  \label{Vergleich_Streams_Signal}
\end{figure}

Neben der DSV-Bibliothek soll es außerdem eine Audio-Bibliothek geben, welche 
auf Audiosignal spezialisiert wird. Der größte Unterschied 
zwischen diesen liegt darin, dass die Audio-Bibliothek das Signal selber erzeugt
und der DSV-Bibliothek ein Signal übergeben wird.

Wegen der Popularität und des niedrigen Preises für die Leistung des ESP32-S3 
Mikrocontroller wird die Signalbibliothek auf diesem entwickelt und
soll neben der eigentlichen 
Signalverarbeitung auch die Anbindung an die benötigte Hardware ermöglichen. Hierzu 
wird eine Trennung zwischen der Signalverarbeitung und hardwarespezifischen Komponenten 
sowie der ESP-IDF
angestrebt. Dadurch soll die Bibliothek plattformunabhängig bleiben. 

Als praktische Anwendung der Bibliothek wird ein digitaler Synthesizer entwickelt.
Der Synthesizer dient dabei sowohl als Anwendung der 
Bibliothek als auch zur praktischen Erprobung ihrer Funktionen.
Dafür wird die nötige Hardware aufgebaut und in der Software implementiert.


% TODO: Schreibe vlt in ein eigenes Kapitel was für tools ich nutze bzw welche ich
% für nützlich halte, auch wenn sie nicht direkt in die Entwicklung der Bibliothek
% eingebunden sind.

% Durunter fallen : Oscilloscope, Logic Analyzer(Kann i2c etc. auslesen),
% VsCode, ESP-IDF, FreRTOS, OpenOCD(on chip debugger), Systemview(Task analyzer), 
% Googletest(für Unit test), OnHardwareTests

% ich denke das mache ich nicht mehr